\documentclass[10pt]{article}

\usepackage[margin=0.75in]{geometry}
\usepackage{amsmath,amsthm,amssymb}
\usepackage{xcolor}
\usepackage{cancel}
\usepackage{graphicx}
\usepackage{changepage}
\usepackage{circuitikz}
\usepackage{pgfplots}
\usepackage{physics}
\usepackage{hyperref}
\usepackage{siunitx}
\usepackage{fontspec}
\usepackage{relsize}
\usepackage{subfig}
\usepackage{todonotes}
\usepackage{multicol, multirow, booktabs}
\usepackage[breakable]{tcolorbox}
\usepackage[inline]{enumitem}

\theoremstyle{definition}
\newtheorem{problem}{Problem}
\newtheorem{soln}{Solution}
\newtheorem{theorem}{Theorem}

\pgfplotsset{compat=newest}
\usetikzlibrary{lindenmayersystems}
\usetikzlibrary{arrows}
\usetikzlibrary{calc}
\usetikzlibrary{positioning, fit}
\usetikzlibrary{3d, perspective}

\definecolor{incolor}{HTML}{303F9F}
\definecolor{outcolor}{HTML}{D84315}
\definecolor{cellborder}{HTML}{CFCFCF}
\definecolor{cellbackground}{HTML}{F7F7F7}
\newcommand{\ui}{\hat{i}}
\newcommand{\uj}{\hat{j}}
\newcommand{\uk}{\hat{k}}
\newcommand{\ux}{\hat{x}}
\newcommand{\uy}{\hat{y}}
\newcommand{\uz}{\hat{z}}
\newcommand{\pr}[1]{{#1}^\prime}
\pgfdeclarelayer{background}  
\pgfsetlayers{background,main}
\AtBeginDocument{\RenewCommandCopy\qty\SI}
\newcommand{\justif}[2]{&{#1}&\text{#2}}

\makeatletter
\newcommand{\boxspacing}{\kern\kvtcb@left@rule\kern\kvtcb@boxsep}
\makeatother
\newcommand{\prompt}[4]{
    \ttfamily\llap{{\color{#2}[#3]:\hspace{3pt}#4}}\vspace{-\baselineskip}
}

\newcommand{\thevenin}[2]{
  \begin{center}
    \begin{circuitikz} \draw
      (0,0) -- (2,0) to[battery1, l_=$V_{Th}\eq#1$] (2,2) 
      to[resistor, l_=$R_{Th}\eq#2$] (0,2)
      ;
      \draw [o-] (-.07,2.079);
      \draw [o-] (-.07,0.079);
    \end{circuitikz}
  \end{center}
}

\newcommand{\norton}[2]{
  \begin{center}
    \begin{circuitikz} \draw
      (0,0) -- (3,0) to[american current source, l_=$I_{N}\eq#1$] (3,2) -- (0,2) (2,0)
      to[resistor, l=$R_{N}\eq#2$] (2,2)
      ;
      \draw [o-] (-.07,2.079);
      \draw [o-] (-.07,0.079);
    \end{circuitikz}
  \end{center}
}

\newcommand{\highlight}[1]{\colorbox{yellow}{$\displaystyle #1$}}

\newcommand{\ti}[1]{\widetilde{#1}}

\newfontface{\Kaufmann}{Kaufmann}
\DeclareTextFontCommand{\kf}{\Kaufmann}
\newcommand{\scriptr}{\fontsize{12pt}{12pt}\kf{r}}

\newfontface{\KaufmannB}{Kaufmann Bd BT}
\DeclareTextFontCommand{\kfb}{\KaufmannB}
\newcommand{\bscriptr}{\fontsize{12pt}{12pt}\kfb{r}}

\newcommand{\bv}[1]{\mathbf{#1}}

\title{Math 4360H: Assignment II}
\author{Jeremy Favro (0805980) \\ Trent University, Peterborough, ON, Canada}
\date{\today}

\begin{document}
\maketitle

% PROBLEM 1
\begin{problem}~
\begin{enumerate}[label=(\alph*)]
  \item Let $F$ be a field and $f(x)\in F[x]$. Suppose for some $\alpha\in F$ that
        $f(x+\alpha)$ is irreducible. Show that $f (x)$ must be irreducible.
  \item Show that $x^4 + 4x + 1$ is irreducible over $\mathbb{Q}$ by using part (a)
        with $\alpha=1$.
\end{enumerate}
\end{problem}
\begin{soln}~
  \begin{enumerate}[label=(\alph*)]
    \item If $f(x+\alpha)$ is irreducible over $F$ then $f(x+\alpha)=g(x+\alpha)h(x+\alpha)$ and one of $g$ or $h$ is constant. We can also write that
          $f(x)=g(x)h(x)$ and since either $g(x+\alpha)$ or $h(x+\alpha)$ must be constant so must one of $g(x)$ or $h(x)$ so $f(x)$ is irreducible over $F$.
    \item $f(x)=x^4 + 4x + 1$, $f(x+1)=(x+1)^4 + 4(x+1) + 1=x^4+4x^3+6x^2+4x+1+4x+4+1=x^4+4x^3+6x^2+8x+6$ which is irreducible over $\mathbb{Q}$
          by Eisenstein's criterion with $p=2$.
  \end{enumerate}
\end{soln}

% PROBLEM 2
\begin{problem}
Show that $f (x) = x^5 +3x^2 +1$ is irreducible over $\mathbb{Q}$ using the method
of Example 9.20 in the notes.
\end{problem}
\begin{soln}
  There are several ways this could factor. First let's eliminate the linear case. By corollary 9.16 the roots of $f(x)$ can be $\pm 1$. Since neither of these is a root
  we cannot factor out a linear factor. So this must factor (if it is reducible) as an irreducible cubic and an irreducible quadratic. So,
  $$f(x)=(ax^3+bx^2+cx+d)(ex^2+fx+g).$$
  Multiplying this out we obtain that
  $$f(x)=x^5 +3x^2 +1=aex^5 + (b e + a f)x^4 + (c e + b f + a g) x^3 + (d e + c f + b g) x^2  + (d f + c g) x + d g$$
  which gives us the system of equations
  \begin{align*}
    ae              & =1 \\
    b e + a f       & =0 \\
    c e + b f + a g & =0 \\
    d e + c f + b g & =3 \\
    d f + c g       & =0 \\
    d g             & =1
  \end{align*}
\end{soln}

% PROBLEM 3
\begin{problem}
\begin{enumerate}[label=(\alph*)]
  \item Show that the map $\phi_n:\mathbb{Z}[x]\to \mathbb{Z}_n[x]$ given by
        $$\phi_n\left(\sum_{i=0}^{k}a_ix^i\right)=\sum_{i=0}^{k}\left(a_i\mod n\right)x^i$$
        is a ring homomorphism.
  \item Let $f (x) = a_0 + a_1x + \dots + a_mx^m\in\mathbb{Z}[x]$, with $a_m\neq 0$, and let
        $p$ be a prime which does not divide $a_m$. Prove that if $\phi_p(f(x))$
        is irreducible over $\mathbb{Z}_p$, then $f (x)$ is irreducible over $\mathbb{Q}$.
  \item Is the converse of part (b) true?
  \item Use part (b) to show that each of the following polynomials are
        irreducible over $\mathbb{Q}$.
        \begin{enumerate}[label=(\roman*)]
          \item $7x^3 + 6x^2 + 4x + 6$
          \item $21x^3 - 3x^2 + 2x + 9$
          \item $x^4 + 7x^3 + 16x^2 + 8x + 11$
          \item $3x^4 + 9x^3 + 6x^2 + 3x + 1$
          \item $x^5 + 2x + 4$
        \end{enumerate}
\end{enumerate}
\end{problem}
\begin{soln}
  \begin{enumerate}[label=(\alph*)]
    \item A ring homomorphism is a map which preserves both addition and multiplication. I'll start with addition,
          Let $\left\{a_i\right\},\left\{b_i\right\}\in \mathbb{Z}$ be sets of coefficients of two polynomials. These polynomials can each be written as
          $$\sum_{i=0}^{k}a_ix^i\quad\text{and}\quad \sum_{i=0}^{k}b_ix^i$$
          even if they differ in degree, provided we take the sets of coefficients to be zero past some point. Here $k$ is the higher degree (``index'' (not sure if that's a thing with sets)) of the last nonzero element,
          so that we get the entirety of the higher degree polynomial. Now,
          \begin{align*}
            \phi_n\left(\sum_{i=0}^{k}a_ix^i+\sum_{i=0}^{k}b_ix^i\right) & =\phi_n\left(\sum_{i=0}^{k}a_ix^i+b_ix^i\right)                                           \\
                                                                         & =\phi_n\left(\sum_{i=0}^{k}\left[a_i+b_i\right]x^i\right)\justif{\quad}{Distributive law} \\
                                                                         & =\sum_{i=0}^{k}\left[\left(a_i+_nb_i\right)\mod n\right] x^i
          \end{align*}
          and we know that $a_i+_nb_i=\left(a_i\mod n+b_i\mod n\right)\mod n$ and $c\mod n \mod n=c\mod n$ so
          \begin{align*}
            \sum_{i=0}^{k}\left[\left(a_i+_nb_i\right)\mod n\right] x^i & =\sum_{i=0}^{k}\left[\left(a_i\mod n+b_i\mod n\right)\mod n\right] x^i           \\
                                                                        & =\sum_{i=0}^{k}\left[\left(a_i\mod n+b_i\mod n \right)x^i\mod n\right]           \\
                                                                        & =\sum_{i=0}^{k}\left(a_i\mod n\right)x^i+\sum_{i=0}^{k}\left(b_i\mod n\right)x^i \\
                                                                        & =\phi_n\left(\sum_{i=0}^{k}a_ix^i\right)+\phi_n\left(\sum_{i=0}^{k}b_ix^i\right)
          \end{align*}
          Now under multiplication,
          \begin{align*}
            \phi_n\left(\sum_{i=0}^{k}a_ix^i\sum_{i=0}^{k}b_ix^i\right) & = \phi_n\left(\sum_{i=0}^{2k}\left[x^i\sum_{m+l=i}a_mb_l\right]\right)                     \\
                                                                        & = \sum_{i=0}^{2k}\left[x^i\sum_{m+l=i}a_m\cdot_n b_l\mod n\right]                          \\
                                                                        & = \sum_{i=0}^{2k}\left[x^i\sum_{m+l=i}\left(a_m\mod n \cdot  b_l\mod n\right)\mod n\right] \\
                                                                        & = \sum_{i=0}^{k}x^ia_i\mod n\sum_{i=0}^{k}b_i\mod n                                        \\
                                                                        & =\phi_n\left(\sum_{i=0}^{k}a_ix^i\right)\phi_n\left(\sum_{i=0}^{k}b_ix^i\right)
          \end{align*}
    \item We are given that $\phi_p(f(x))$ is irreducible. Suppose that $f(x)$ is reducible and hence $f(x)=g(x)h(x)$ for some non-constant $g(x),h(x)\in\mathbb{Z}[x]$.
          Then since $\phi_n$ is a homomorphism, $$\phi_p(f(x))=\phi_p(g(x)h(x))=\phi_p(g(x))\phi_p(h(x))$$
          and since $\phi_p(f(x))$ is irreducible, one of $\phi_p(g(x)),\phi_p(h(x))$ must be constant and hence $p$ must divide all its coefficients except the constant one
          (assume WLOG that this is $h(x)$).
          If we write that
          $g(x)=g_0+\dots+g_gx^g$ (where I'm representing $\deg(g(x))=g$) and $h(x)=h_0+\dots+h_hx^h$ then
          $$f(x)=\sum_{i=0}^{g+h}x^i\left[\sum_{j+k=i}g_jh_k\right]$$ but we know that $p\mid h_k\forall k>0$ and that $p\nmid a_m$, the highest order coefficient of $f(x)$, however $a_m$ contains
          a coefficient of $h(x)$ inside it which \emph{is} divided by $p$, hence we have a contradiction and $f(x)$ is irreducible over $\mathbb{Z}$. In class we proved that
          if $f (x)$ is reducible over $\mathbb{Q}$, then it is reducible over $\mathbb{Z}$. Taking the contrapositive of this we obtain that
          if $f (x)$ is irreducible over $\mathbb{Z}$, then it is irreducible over $\mathbb{Q}$, and we have proved here that $\phi_p(f(x))$ being irreducible over $\mathbb{Z}_p$
          implies that $f(x)$ is irreducible over $\mathbb{Z}$ and hence by the contrapositive of the theorem, it is irreducible over $\mathbb{Q}$.
    \item The converse would be that if $f(x)$ is irreducible over $\mathbb{Q}$ then $\phi_p(f(x))$ is irreducible over $\mathbb{Z}_p$. This is not generally true. Take
          $x^2-2$ in $\mathbb{Z}_7$ where $x^2-2=(x-3)(x+3)$, however $x^2-2$ has no roots in $\mathbb{Q}$ and hence the contrapositive is false and hence the original statement is false.
    \item \begin{enumerate}[label=(\roman*)]
            \item Here we take $p=5$ so that $\phi_5(7x^3 + 6x^2 + 4x + 6)=2x^3 + x^2 + 4x + 1$ which has no roots in $\mathbb{Z}_5$ and hence does not reduce there and hence does not reduce in $\mathbb{Q}$.
            \item Here we take $p=3$ so that $\phi_3(21x^3 - 3x^2 + 2x + 9)=2x$ which clearly is irreducible over anything.
            \item Here we take $p=2$ so that $\phi_2(x^4 + 7x^3 + 16x^2 + 8x + 11)=x^4 + x^3 + 1$ which is irreducible over $\mathbb{Z}_2$ as it has no roots and hence is irreducible over $\mathbb{Q}$.
            \item Here we take $p=3$ so that $\phi_3(3x^4 + 9x^3 + 6x^2 + 3x + 1)=1$ which does not reduce anywhere.
            \item Here we take $p=3$ so that $\phi_3(x^5 + 2x + 4)=x^5+2x+1$ which is irreducible over $\mathbb{Z}_3$ as it has no roots and hence is irreducible over $\mathbb{Q}$.
          \end{enumerate}
  \end{enumerate}
\end{soln}

% PROBLEM 4
\begin{problem}
Let F be a field. Consider the set
$$\ti{F}[x]=\left\{a_0+a_1x+a_2x^2+\dots+a_nx^n\in F[x]|a_0\neq 0\right\}$$
\begin{enumerate}[label=(\alph*)]
  \item Show that $\ti{F}[x]$ is closed under multiplication.
  \item Let $f(x)\in \ti{F}[x]$, Show that if $g(x)$ divides $f(x)$ then $g(x)\in \ti{F}[x]$.
  \item Consider the map $\rho: \ti{F}[x]\to \ti{F}[x]$ given by
        $$\rho(a_0+a_1x+a_2x^2+\dots+a_nx^n)=a_n+a_{n-1}x+a_{n-2}x^2+\dots+a_0x^n$$
        where $a_n\neq 0$. Show that $$\rho(f(x)g(x))=\rho(f(x))\rho(g(x))$$
        for $f(x),g(x)\in \ti{F}[x]$.
  \item Let $f(x)\in\ti{F}[x]$. Show that if $\rho(f(x))$ is irreducible over $F$, then $f(x)$ is irreducible over $F$.
  \item Verify part (d) for $f(x)=x^3+2x^2+3\in \mathbb{Z}_5[x]$.
  \item Show that the polynomial $3x^4 + 9x^3 + 6x^2 + 3x + 1$ is irreducible
        over $\mathbb{Q}$.
  \item Show that the polynomial $px^n -1$  is irreducible over $\mathbb{Q}$ for any
        prime $p$ and $n\in\mathbb{N}$. Give an example of $a\in\mathbb{Z}$ (not prime) and
        $n\in\mathbb{N}$ such that $ax^n - 1$ is reducible over $\mathbb{Q}$.
\end{enumerate}
\end{problem}
\begin{soln}
  \begin{enumerate}[label=(\alph*)]
    \item The $j$th coefficient of the product of two polynomials with coefficients $\left\{a_i\right\}$ and $\left\{b_i\right\}$ is given by
          $$\sum_{m+n=j}a_mb_n.$$
          Since $j=0$ for the constant term, it is given by $a_0b_0$. Since these are elements of a field, there are no zero divisors and hence $a_0b_0\neq 0$, so
          the product of two polynomials with nonzero constant term will always itself have nonzero constant term (over a field) and hence will belong to $\ti{F}[x]$, meaning it is closed under multiplication.
    \item If $g(x)$ divides $f(x)$ then $f(x)=q(x)g(x)$. By part (a) this is only possible if the product of the constant terms of $q(x)$ and $g(x)$ is nonzero
          and hence $g(x)$ must have nonzero constant term, making it an element of $\ti{F}[x]$. Note that I am assuming that $g(x)$ comes from a structure in which there are non zero divisors.
    \item What $\rho$ does is map the coefficient of $x^j$ to the coefficient of $x^{n-j}$ where $n$ is the polynomial's degree. Hence for the product of two polynomials with coefficients $\left\{a_i\right\}$, $\left\{b_i\right\}$
          whose $j$th coefficient is given by $$\sum_{j=k+l}a_kb_l$$ $\rho$ will again map $j$ to $n-j$ so that the new coefficient of $x^j$ in $\rho(f(x)g(x))$ is given by
          $$\sum_{n-j=k+l}a_kb_l.$$
          Note that here $n$ is the degree of the larger polynomial.
          And the coefficients product
          $\rho(f(x))\rho(g(x))$ is given by the usual summation, except the indices of the coefficients must now add to $n-j$, as this is the new power of $x$ that they belong to, so:
          $$\sum_{n-j=k+l}a_kb_l$$
          for the coefficients in $\rho(f(x))\rho(g(x))$. Hence since polynomials are compared by coefficients, $\rho(f(x)g(x))=\rho(f(x))\rho(g(x))$.
    \item Suppose $f(x)$ is reducible over $F$. Then $f(x)=g(x)h(x)$ and $\rho(f(x))=\rho(g(x)h(x))=\rho(g(x))\rho(h(x))$. Since we are assuming that $\rho(f(x))$ is irreducible,
          one of $\rho(h(x))$ or $\rho(g(x))$ must be constant (assume WLOG it's $h(x)$). Because this is the case then $h(x)$ itself must be constant by the definition of $\rho$ and hence
          $f(x)$ is irreducible as it only factors into a polynomial times a constant.
    \item Since $\rho(x^3+2x^2+3)=3x^3+2x+1$ and this is irreducible since it has no zeroes in $\mathbb{Z}_5$.
    \item Since $\rho(3x^4 + 9x^3 + 6x^2 + 3x + 1)=x^4 + 3x^3 + 6x^2 + 9x + 3$ for which there exists $p=3$ which satisfies Eisenstein's criterion, meaning $\rho(f(x))$
          is irreducible over $\mathbb{Q}$ and so is $f(x)$.
    \item Since $\rho(px^n -1)=p-x^n$ which is irreducible over $\mathbb{Q}$ by Eisenstein's criterion with $p=p$, $px^n -1$ is also irreducible over $\mathbb{Q}$.
          $x^2-1$ is an example where $a=1$ that is reducible over $\mathbb{Q}$. It reduces to $(x-1)(x+1)$.
  \end{enumerate}
\end{soln}

% PROBLEM 5
\begin{problem}
The following is known as Brauer's Irreducibility Criterion.
\begin{theorem}
  Let $a_0\leq a_1\leq\dots\leq a_{n-1}$ be positive integers and $n\geq 2$. Then the polynomial $f(x)=x^n-a_{n-1}x^{n-1}-\dots-a_1x-a_0$ is irreducible over $\mathbb{Q}$.
\end{theorem}
Let $c$ be a nonzero integer, and let $n\in\mathbb{N}$ with $n \geq 3$. Show that
the polynomial $f(x)=x^n+cx^{n-1}+cx^{n-2}+\dots+cx-1$ is irreducible
over $\mathbb{Q}$.
\end{problem}
\begin{soln}
Okay, taking the fact that I'm give Brauer's Criterion as a hint I'm going to try and make $f(x)=x^n+cx^{n-1}+cx^{n-2}+\dots+cx-1$ look like the polynomial in the theorem.
We can do this by taking \begin{equation}
  -x^nf(\frac{1}{x})=-1-cx-cx^{2}-\dots-cx^{n-1}+x^n.
\end{equation}
But where to from here? Well in the case that $c>0$ the action of taking $x^nf(\frac{1}{x})$ is the same as $\rho(f(x))$ from problem 4 and since (1) is irreducible by 
Brauer's Criterion $f(x)$ is also irreducible by the properties we proved in question 4. For $c<0$ if we just take 
\begin{equation}
  x^nf(\frac{1}{x})=1+cx+cx^{2}+\dots+cx^{n-1}-x^n=1-kx-kx^{2}-\dots-kx^{n-1}-x^n
\end{equation}
for $k=\abs{c}$. If we take $\rho$ of this we get 
$$-1-kx-kx^{2}-\dots-kx^{n-1}+x^n$$
which is irreducible by Brauer's criterion, and hence $x^nf(\frac{1}{x})$ is irreducible by question 4 and so $f(x)$ is irreducible by question 4.
\end{soln}
\end{document}